\subsection{Effective descent data}

Let X and Y be topological spaces, \(f: X \to Y\) a continuous map, and let \((Z, p)\) be an X-space. Let \(\tau\) be a descent datum relative to f on \((Z, p)\) , let \(R_{\tau}\) be the associated equivalence relation and let \(g: Z \to Z/R_{\tau}\) be the canonical map. Since the relation \(R_{\tau}\) is compatible with the map \(f \circ p\) , there exists a unique continuous map \(q: Z/R_{\tau} \to Y\) such that the diagram

\[
\begin{array}{c} \mathrm{Z} \xrightarrow {g} \mathrm{Z} / \mathrm{R} _ {\tau} \\ p \Biggl \downarrow \qquad \qquad \qquad \qquad \qquad \qquad \qquad \Biggl \downarrow q \\ \mathrm{X} \xrightarrow {f} \mathrm{Y} \end{array}\tag{1}
\]

The Y-space \((\mathrm{Z}/\mathrm{R}_{\tau}, q)\) is called the quotient space of \((\mathrm{Z}, p)\) by the descent datum \(\tau\) . Denote by \(h: Z \to X \times_{Y} (Z/R_{\tau})\) the map defined by \(h(z) = (p(z), g(z))\) for \(z \in Z\) . It is continuous. Let \((x, u) \in X \times_{Y} (Z/R_{\tau})\) and let \(z \in Z\) such that \(g(z) = u\) ; we have \((z, x) \in Z \times_{Y} X\) and the point \(z' = \tau(z, x)\) is the unique element of Z such that \(h(z') = (x, u)\) ; consequently, the map h is bijective.

One says that the descent datum \(\tau\) relative to f on \((\mathrm{Z},p)\) is effective if diagram (1) is a Cartesian square, that is, if the continuous bijection h is a homeomorphism. For the descent datum \(\tau\) to be effective, it is necessary and sufficient that the sets \(p^{-1}(\mathrm{U})\cap \mathrm{V}\) , where U is an open subset of X and V an open subset of Z saturated for \(R_{\tau}\) , form a basis of the topology of Z. In particular, the condition, for a descent datum relative to f, of being effective is local in nature in Y.

Examples. --- 1) Let X and Y be topological spaces and \(f\colon X\to Y\) be a continuous map. Let \((T,q)\) be a Y-space. Let Z be the X-space \(X\times_Y T\) , equipped with the map \(\mathrm{pr}_{1}\) ; let us denote by \(\tau\) its canonical descent datum relative to f (IV, p.~383, example 1). The subsets of Z of the form \(X\times_Y V\) where V is an open subset of T, are open in Z and saturated for the relation \(\mathrm{R}_{\tau}\) (loc. cit.). By definition of the product topology, the sets \(U\times_Y\overline{\mathrm{pr}}_{2}^{-1}(V)\) , where U is open in X and V is open in T, form a basis for the topology of \(X\times_Y Z\) . Consequently, the canonical descent datum on a fiber product \(X\times_Y T\) is effective.

The canonical map \(Z/R_{\tau} \rightarrow T\) is injective and continuous. It is not, however, necessarily surjective, nor strict (IV, p.~462, exerc. 1).

\begin{enumerate}
\def\labelenumi{\arabic{enumi})}
\setcounter{enumi}{1}
\tightlist
\item
  Let us resume the notation of Example 2 (IV, p.~383). The topological space \(Z/R_{\tau}\) is then the topological space obtained by gluing the spaces \(Z_{i}\) along the \(p_{i}^{-1}(V_{i} \cap V_{j})\) by means of the bijections \(\tau_{i,j}\) . Hence, if for every \(i \in I\) the set \(V_{i}\) is open (resp. closed) in Y, the set \(g(Z_{i})\) is open (resp. closed) in \(Z/R_{\tau}\) and the restriction of g to \(Z_{i}\) induces a homeomorphism of \(Z_{i}\) onto \(g(Z_{i})\) (TG, I, p.~17, prop. 9). The space Z is the sum space of the spaces \(Z_{i}\) ; the space \(X \times_{Y}(Z/R_{\tau})\) is the direct sum space of the spaces \(V_{i} \times_{Y}(Z/R_{\tau}) = g(Z_{i})\) . The map h is identified with the sum map of the maps \(g|Z_{i}: Z_{i} \to g(Z_{i})\) . It is therefore a homeomorphism, which proves that the descent datum \(\tau\) is effective.
\end{enumerate}

Without any particular assumption on the sets \(V_{i}\) it is not always true that the restriction of g to \(Z_{i}\) induces a homeomorphism of \(Z_{i}\) onto its image; in this case, the descent datum \(\tau\) is not effective (IV, p.~462, exerc. 2).

PROPOSITION 1. --- Let X and Y be topological spaces, let \(f: X \to Y\) be a continuous map. Suppose that every point of Y has a neighborhood over which there exists a continuous section of the map f . Every descent datum relative to f on an X-space is then effective.

Let \((Z,p)\) be an X-space and let \(\tau\) be a descent datum relative to f on \((Z,p)\) . The assertion that \(\tau\) is an effective descent datum is local in Y, which allows one to assume that the map f has a continuous section s. Denote by \(g: Z \to Z/R_{\tau}\) and \(q: Z/R_{\tau} \to Y\) the canonical maps and let \(h: Z \to X \times_{Y} (Z/R_{\tau})\) be the map given by \(z \mapsto (p(z), g(z))\) . The map h is bijective and continuous; it suffices to show that it is a homeomorphism.

The map from Z to Z which sends z to \(\tau(z,s(f(p(z))))\) is continuous and sends every element of Z to the unique element \(z'\) of Z equivalent to it for the relation \(R_{\tau}\) and such that \(p(z')\) belongs to the image of s. It therefore defines, by passage to the quotient, a continuous map \(t:Z/R_{\tau}\to Z\) which is a section of the map g. In particular, one has \(f\circ p\circ t=q\circ g\circ t=q\) .

For every \((x,u)\in\mathrm{X}\times_{\mathrm{Y}}(\mathrm{Z}/\mathrm{R}_{\tau})\) , we have \(f(p(t(u)))=f(x)\) ; then set \(h'(x,u)=\tau(t(u),x)\) . The map \(h'\) from \(X\times_{Y}(Z/R_{\tau})\) into Z thus defined is continuous. For every \((x,u)\in\mathrm{X}\times_{\mathrm{Y}}(\mathrm{Z}/\mathrm{R}_{\tau})\) , we have

\[
\begin{array}{r l} h (h ^ {\prime} (x, u)) & = (p (h ^ {\prime} (x, u)), g (h ^ {\prime} (x, u))) \\ & = (p (\tau (t (u), x)), g (\tau (t (u), x))) \\ & = (x, u), \end{array}
\]

since \(\tau(t(u), x)\) is equivalent to \(t(u)\) for the relation \(\mathrm{R}_{\tau}\) . This proves that the map \(h \circ h'\) is the identity map of \(\mathrm{X} \times_{\mathrm{Y}} (\mathrm{Z}/\mathrm{R}_{\tau})\) . For \(z \in \mathrm{Z}\) , we then have \(t(g(z)) = \tau(z, s(f(p(z))))\) and \(f(p(t(g(z)))) = f(p(z))\) , whence \(z = \tau(t(g(z)), p(z))\) , by definition of the equivalence relation \(\mathrm{R}_{\tau}\) . We therefore have \(h'(h(z)) = z\) and \(h' \circ h = \mathrm{Id}_{\mathrm{Z}}\) . The map \(h\) is therefore a homeomorphism, which is what had to be proved.

PROPOSITION 2. --- Let \(f\colon X \to Y\) be a continuous map, let \((Z,p)\) be an X-space and let \(\tau\) be a descent datum relative to f on \((Z,p)\) . The equivalence relation \(R_{\tau}\) is closed if f is proper; it is open if f is open.

The map \(\widetilde{\tau}\colon\mathbf{Z}\times_{\mathrm{Y}}\mathbf{X}\to\mathbf{X}\times_{\mathrm{Y}}\mathbf{Z}\) given by \((z,x)\mapsto(p(z),\tau(z,x))\) is a homeomorphism, with inverse mapping \((x,z)\mapsto(\tau(z,x),p(z)).\) We have \(\tau=\mathrm{pr}_{2}\circ\widetilde{\tau}\) , where \(\mathrm{pr}_{2}\colon\mathbf{X}\times_{\mathrm{Y}}\mathbf{Z}\to\mathbf{Z}\) is the second projection. If f is proper, \(\mathrm{pr}_{2}\) is proper; if f is open, \(\mathrm{pr}_{2}\) is open (I, p.~17, prop. 8). It follows that \(\tau\) is proper (resp. open) if f is. The saturation of a subset A of Z with respect to the relation \(\mathrm{R}_{\tau}\) is the image of \(\mathbf{A}\times_{\mathrm{Y}}\mathbf{X}\) under \(\tau\) . Consequently, if f is proper, the saturation of a closed subset is closed; if f is open, the saturation of an open subset is open.

